% Source: Weinberg GR, physical PDF pp. 382--384
%         (printed pp. 361--363).
% Coverage: Section 12.3, General Covariance and Energy-Momentum
%           Conservation; equations (12.3.1)--(12.3.2).
% Figures/tables/footnotes: none.
% Status: source-reviewed and compile-clean.
% Uncertainties: coordinate displacements, Lie derivatives, and compact
%                conservation laws are documented below and source-checked.
%
% MODERNIZATION CHECK: Weinberg's coordinate displacement epsilon^mu is xi^mu
% here, and the field changes in (12.3.1) are identified explicitly as minus
% Lie derivatives. Comma derivatives in (12.3.2) and in the current
% conservation law are replaced by nabla.
% LITERAL-OPERATOR AUDIT: the minus sign in front of the metric variation,
% the minus metric-derivative term after integration by parts, and the minus
% sign in the integrated gauge variation were checked against pp. 383--384.

\subsection{General Covariance and Energy-Momentum Conservation}
\label{sec:12.3}

If the action \(I_M\) for a material system is a scalar, then the statement
that \(\delta I_M\) vanishes is generally covariant, and so also are the
dynamical equations derived from this statement. For instance, a glance at
the action \eqref{eq:12.1.6} for a collisionless plasma shows that this
\(I_M\) is a scalar, and this ensures the general covariance of the dynamical
equations \eqref{eq:12.1.1}--\eqref{eq:12.1.3} that follow from
Eq.~\eqref{eq:12.1.6} upon application of the principle of stationary action.

We shall therefore assume that \(I_M\) is a scalar. This means that \(I_M\)
will be unchanged if we simultaneously make the replacements
\[
\begin{gathered}
  \dd^4x\longrightarrow\dd^4x',
  \qquad
  \partial_\mu
  \longrightarrow
  \partial'_\mu
  =
  \frac{\partial x^\nu}{\partial x'^\mu}\partial_\nu,
  \qquad
  x_n^\mu(p)\longrightarrow x_n'^\mu(p),
  \\
  A^\mu(x)
  \longrightarrow
  A'^\mu(x')
  \equiv
  A^\nu(x)\frac{\partial x'^\mu}{\partial x^\nu},
  \\
  g_{\mu\nu}(x)
  \longrightarrow
  g'_{\mu\nu}(x')
  \equiv
  g_{\rho\sigma}(x)
  \frac{\partial x^\rho}{\partial x'^\mu}
  \frac{\partial x^\sigma}{\partial x'^\nu}.
\end{gathered}
\]
However, \(x'^\mu\) is a mere variable of integration (as opposed to
\(x_n'^\mu\), which is a dynamical variable), so we can change \(x'^\mu\)
back to \(x^\mu\) without changing \(I_M\). We conclude then that \(I_M\) is
unchanged by the replacements
\[
\begin{gathered}
  x_n^\mu(p)\longrightarrow x_n'^\mu(p),
  \\
  A_\mu(x)
  \longrightarrow
  A'_\mu(x)
  =
  A_\nu(x)\frac{\partial x^\nu}{\partial x'^\mu}
  -\left[A'_\mu(x')-A'_\mu(x)\right],
  \\
  g_{\mu\nu}(x)
  \longrightarrow
  g'_{\mu\nu}(x)
  =
  g_{\rho\sigma}(x)
  \frac{\partial x^\rho}{\partial x'^\mu}
  \frac{\partial x^\sigma}{\partial x'^\nu}
  -\left[g'_{\mu\nu}(x')-g'_{\mu\nu}(x)\right],
\end{gathered}
\]
with \(\dd^4x\) and \(\partial_\mu\) now left alone. If the original
transformation \(x^\mu\longrightarrow x'^\mu\) was infinitesimal,
\[
  x'^\mu=x^\mu+\xi^\mu(x),
\]
then the changes in the dynamical variables are
\begin{equation}
\begin{aligned}
  \delta x_n^\mu(p)
  &=
  \xi^\mu(x_n(p)),
  \\
  \delta A_\mu(x)
  &=
  -A_\lambda(x)\partial_\mu\xi^\lambda(x)
  -\partial_\lambda A_\mu(x)\,\xi^\lambda(x)
  =
  -(\Lie_\xi A)_\mu,
  \\
  \delta g_{\mu\nu}(x)
  &=
  -g_{\mu\lambda}(x)\partial_\nu\xi^\lambda(x)
  -g_{\lambda\nu}(x)\partial_\mu\xi^\lambda(x)
  -\partial_\lambda g_{\mu\nu}(x)\,\xi^\lambda(x)
  =
  -(\Lie_\xi g)_{\mu\nu}.
\end{aligned}
  \tag{12.3.1}\label{eq:12.3.1}
\end{equation}
(This change in \(A_\mu\) or \(g_{\mu\nu}\) is just the Lie derivative; see
Section~\ref{sec:10.9}.) The important point is that this is now an
infinitesimal transformation of the dynamical variables alone, not of the
coordinates over which we integrate, so the principle of stationary action
tells us that when the dynamical equations for \(x_n^\mu\), \(A_\mu\), and so
on, are satisfied, the change in these quantities produces no change in the
matter action \(I_M\). The only change in \(I_M\) comes from the variation in
the external field \(g_{\mu\nu}\), and Eq.~\eqref{eq:12.2.2} gives for this
change
\[
  \delta I_M
  =
  -\frac12\int\dd^4x\,\sqrt{-g}\,T^{\mu\nu}
  \left[
    g_{\mu\lambda}\partial_\nu\xi^\lambda
    +g_{\lambda\nu}\partial_\mu\xi^\lambda
    +(\partial_\lambda g_{\mu\nu})\xi^\lambda
  \right].
\]
If \(I_M\) is a scalar, then this must vanish; integrating by parts gives
\[
  0=\delta I_M
  =
  \int\dd^4x\,\xi^\lambda
  \left[
    \partial_\nu\!\left(\sqrt{-g}\,T^\nu{}_\lambda\right)
    -\frac12(\partial_\lambda g_{\mu\nu})
     \sqrt{-g}\,T^{\mu\nu}
  \right],
\]
and since \(\xi^\lambda(x)\) is arbitrary,
\[
  0
  =
  \partial_\nu\!\left(\sqrt{-g}\,T^\nu{}_\lambda\right)
  -\frac12(\partial_\lambda g_{\mu\nu})
   \sqrt{-g}\,T^{\mu\nu},
\]
or, recalling Eq.~\eqref{eq:4.7.6},
\begin{equation}
  \nabla_\nu T^\nu{}_\lambda=0.
  \tag{12.3.2}\label{eq:12.3.2}
\end{equation}
\emph{Thus the energy-momentum tensor defined by Eq.~\eqref{eq:12.2.2} is
conserved (in the covariant sense) if and only if the matter action is a
scalar.} Also, with \(I_M\) a scalar, Eq.~\eqref{eq:12.2.2} shows immediately
that \(T^{\mu\nu}\) is a symmetric tensor, so this definition of the
energy-momentum tensor has all the properties for which one could wish.

This proof, that general covariance implies energy-momentum conservation, has
an exact analog in the proof that gauge invariance implies charge
conservation. The change in the action \(I'_M\) defined by
Eq.~\eqref{eq:12.2.3}, caused by an arbitrary gauge transformation, can arise
only from the change in \(A_\mu\), since \(I'_M\) is stationary with respect
to all other dynamical variables. A general infinitesimal gauge
transformation \(\chi\) will produce in \(A_\mu\) the change (see
Sections~\ref{sec:4.11} and \ref{sec:10.2})
\[
  \delta A_\mu=\partial_\mu\chi.
\]
Using this in Eq.~\eqref{eq:12.2.6}, we see that \(I'_M\) is gauge invariant
if and only if
\[
  0=\delta I'_M
  =
  \int\dd^4x\,\sqrt{-g}\,J^\mu\partial_\mu\chi.
\]
Integrating by parts gives
\[
  0
  =
  -\int\dd^4x\,\chi\,
  \partial_\mu\!\left(\sqrt{-g}\,J^\mu\right),
\]
or, since \(\chi\) is arbitrary,
\[
  0
  =
  \frac1{\sqrt{-g}}
  \partial_\mu\!\left(\sqrt{-g}\,J^\mu\right)
  =
  \nabla_\mu J^\mu.
\]
We see again how closely analogous are gauge invariance and general
covariance.
